\chapter{Discussion}
\label{cha:discussion}

\section{Predicting with PCPG}

\subsection{Bias and Variance}
\label{sec:bias-variance}

We compare the bias and variance of MC with PCPG-MC, which samples both expectations, and PCPG, which computes both by quadrature.

\subsubsection{MC Estimator}

The MC estimator draws $N$ samples ${\xm}^{(i)} \sim p(\xm \mid \xo)$ and returns
\[
  \widehat{p}_{\mathrm{MC}}(N) = \frac{1}{N}\sum_{i=1}^{N} \sigma\!\left(w^T x^{(i)}\right).
\]
\paragraph{Bias.}
MC is unbiased for all $N \geq 1$: sampling from $p(\xm \mid \xo)$ and averaging $\sigma(w^T x)$ is a direct Monte Carlo estimate of the target expectation, so $\E[\widehat{p}_{\mathrm{MC}}(N)] = p(y \mid \xo)$ exactly.

\paragraph{Variance.}
The draws are i.i.d., so the variance of their average is the per-sample variance
$\sigma^2_{\mathrm{MC}} := \mathrm{Var}_{p(\xm\mid\xo)}\!\left[\sigma(w^T x)\right]$
divided by $N$. Since $\sigma(\cdot) \in (0,1)$, the per-sample variance is bounded:
for any $[0,1]$-valued random variable $Y$,
\[
  \mathrm{Var}(Y) = \E[Y^2] - \E[Y]^2 \leq \E[Y] - \E[Y]^2 = \E[Y](1-\E[Y]) \leq \tfrac{1}{4},
\]
using $Y^2 \leq Y$ on $[0,1]$ and $p(1-p) \leq 1/4$. Therefore
\[
  \mathrm{Var}\!\left[\widehat{p}_{\mathrm{MC}}(N)\right] = \frac{\sigma^2_{\mathrm{MC}}}{N} \leq \frac{1}{4N}.
\]

\subsubsection{PCPG-MC Estimator}

PCPG-MC draws $M$ outer samples $\gamma_j \sim \PG(1,0)$ and, for each of them, $L$ inner samples $u_{ij} \sim \mathcal{N}(0, \gamma_j)$ (Section~\ref{subsec:pcpg-mc}). With $Z := \tfrac{1}{2}\mathrm{Re}[f(u, \gamma)]$ for a single sample, it returns
\[
  \widehat{p}_{\mathrm{PCPG\text{-}MC}}(M,L) = \frac{1}{ML}\sum_{j=1}^{M}\sum_{i=1}^{L} Z_{ij}.
\]

\paragraph{Bias.}
PCPG-MC is unbiased as well. Both expectations are approximated by ordinary MC, so by linearity of expectation the estimator equals the true double integral, which is $p(y \mid \xo)$ by the PG and Hubbard--Stratonovich identities.

\paragraph{Output range.}
Since $|\phi_{\xm\mid\xo}| \leq 1$, a single sample satisfies $|Z| \leq \tfrac{1}{2}e^{1/(8\gamma)}$.
Because $\gamma \sim \PG(1,0)$ can be arbitrarily small, this bound is unbounded, so individual PCPG samples are not confined to $[0,1]$ and can be negative or exceed $1$, while every MC sample lies in $(0,1)$.

\paragraph{Variance.}
The estimator has two nested levels of randomness. For a fixed $\gamma_j$, the average of its $L$ inner samples has the conditional variance $\sigma^2_{A,\gamma_j}/L$, with $\sigma^2_{A,\gamma} := \mathrm{Var}_{u \mid \gamma}[Z \mid \gamma]$. The law of total variance over both levels and the $M$ independent outer samples give
\[
  \mathrm{Var}\!\left[\widehat{p}_{\mathrm{PCPG\text{-}MC}}(M,L)\right] = \frac{\sigma^2_A}{ML} + \frac{\sigma^2_B}{M},
  \qquad
  \sigma^2_A := \E_\gamma\bigl[\sigma^2_{A,\gamma}\bigr],
  \quad
  \sigma^2_B := \mathrm{Var}_\gamma\!\left[\E_{u \mid \gamma}[Z \mid \gamma]\right].
\]
$\sigma^2_A$ measures how much the inner expectation fluctuates for a typical $\gamma$, and $\sigma^2_B$ how much $\gamma$ itself shifts the result.

\paragraph{The inner variance is infinite.}
The problem comes from completing the square in Trick~2, which introduced the factor $e^{1/(8\gamma)}$. For a fixed $z$, the inner expectation is known exactly from the HS identity,
\[
  \E_{u \sim \mathcal{N}(0,\gamma)}\!\left[ e^{-iu(z - c)} \right] = \exp\!\left( -\tfrac{1}{2}\gamma \left( z - c \right)^2 \right) \approx e^{-1/(8\gamma)} \quad \text{for small } \gamma,
\]
which is tiny. Multiplied by $e^{1/(8\gamma)}$ it gives a value of order one. So the inner average has to produce a tiny number from terms of modulus one that almost cancel, and every bit of noise in that average is amplified by $e^{1/(8\gamma)}$, which is already about $2.7 \cdot 10^{5}$ for $\gamma = 0.01$.
In terms of variance, a single inner sample has the second moment $e^{1/(4\gamma)}$. Near zero the density of $\PG(1,0)$ decays only like $e^{-1/(8\gamma)}$, which is not fast enough to compensate, so averaged over $\gamma$ the second moment is infinite: formally $\sigma^2_A = \infty$.
Small values of $\gamma$ are rare, which is why the estimator still works in practice, but when they occur they produce extreme samples.

\paragraph{The outer variance shrinks with more missing features.}
For a GMM with diagonal covariances, the characteristic function of the missing features satisfies
\[
  \bigl|\phi_{\xm\mid\xo}(-uw_m)\bigr| \leq \exp\!\left(-\frac{u^2}{2}\,D_m\right), \qquad D_m := \min_k \sum_{j \in \mathrm{mis}} w_j^2\,\sigma_{kj}^2,
\]
where $D_m$ is the variance of the missing part of $w^T x$. When $D_m$ is large, the integrand is nearly zero unless $|u|$ is small, so the inner expectation is close to zero for a wide range of $\gamma$ and hardly varies with it: $\sigma^2_B \approx 0$. More missing features increase $D_m$, so PCPG-MC gains relative to MC as the missing rate grows, while MC behaves the opposite way, since more missing features make $w^T x$ more random and increase $\sigma^2_{\mathrm{MC}}$.

\subsubsection{PCPG Estimator}

\paragraph{Variance.}
Replacing both expectations with quadrature (Gauss--PG for the outer, Gauss--Hermite for the inner) makes the estimator fully deterministic:
\[
  \mathrm{Var}\!\left[\widehat{p}_{\mathrm{PCPG}}\right] = 0.
\]
Its output is bounded, since the nodes $\gamma_j$ are fixed positive numbers and the factor $e^{1/(8\gamma_j)}$ is a known finite constant for each of them. The amplification is mild, because the smallest Gauss--PG node is not close to zero ($\gamma_{\min} \approx 0.044$ for 20 nodes, giving $e^{1/(8\gamma_{\min})} \approx 17$).

\paragraph{Bias.}
Both expectations introduce a quadrature error. Gauss--Hermite with $N$ nodes is exact for polynomials in $u$ up to degree $2N - 1$, and for analytic integrands its error decays super-algebraically in $N$. The Gauss--PG rule with $M$ nodes is exact for polynomials in $\gamma$ up to degree $2M-1$. At $M = 20$ and $N = 60$, both errors are numerically negligible for moderate $|w^T x|$. For large positive $w^T x$ the outer error is amplified by $e^{z/2}$, which the sign flip (Section~\ref{subsec:sign-flip}) removes; with the flip, the outer rules fitted directly to the sigmoid (Appendix~\ref{subsubsec:sigmoid-fit}) are more accurate than Gauss--PG with fewer nodes. The inner Gauss--Hermite error for large $|w^T x|$ is not affected by the flip.
The infinite inner variance of PCPG-MC does not disappear completely, however; it turns into quadrature error, which grows with the scale of $z = w^T x$. Section~\ref{sec:accuracy-pcpg} discusses when this error and the error of the outer rule become significant.

\subsubsection{Comparison}

Table~\ref{tab:estimator-comparison} summarises the properties of the three
estimators.

\begin{table}[ht]
\centering
\caption{Summary of the three estimators. MC uses $N$ samples, PCPG-MC $M$ outer
and $L$ inner samples, PCPG $M$ Gauss--PG and $N$ Gauss--Hermite nodes.
``CF'' means the estimator requires evaluating $\phi_{\xm\mid\xo}$
rather than drawing samples from $p(\xm \mid \xo)$.}
\label{tab:estimator-comparison}
\renewcommand{\arraystretch}{1.4}
\begin{tabular}{llll}
\hline
\textbf{Estimator} & \textbf{Variance} & \textbf{Bias} & \textbf{Needs} \\
\hline
MC      & $\sigma^2_{\mathrm{MC}}/N \leq 1/(4N)$ & $0$ & samples \\
PCPG-MC & $\sigma^2_A/(ML) + \sigma^2_B/M$, with $\sigma^2_A = \infty$ & $0$ & CF \\
PCPG    & $0$ & quadrature error & CF \\
\hline
\end{tabular}
\end{table}

\paragraph{Bias.}
MC and PCPG-MC are exactly unbiased. PCPG introduces a quadrature error, which
decays fast in the number of nodes and is negligible at the parameter choices
used ($M = 20$ Gauss--PG and $N = 60$ Gauss--Hermite nodes), as long as
$|w^T x|$ is moderate.

\paragraph{Variance.}
MC has the variance $\sigma^2_{\mathrm{MC}}/N$, bounded by $1/(4N)$. PCPG-MC is
unbiased as well, but its inner variance $\sigma^2_A$ is formally infinite because
of the factor $e^{1/(8\gamma)}$, so rare small values of $\gamma$ produce extreme
samples, and its outer variance $\sigma^2_B$ only becomes small when many features
are missing. PCPG has no variance at all: its error is pure bias from the
quadrature.

\paragraph{When to use which.}
MC is the simplest choice when sampling from $p(\xm \mid \xo)$ is easy. PCPG-MC
mainly serves as a reference: it shows that the identities behind PCPG are exact,
but sampling the two expectations is not competitive. PCPG is the method we
propose: it only needs the characteristic function, and it is deterministic, so
repeated evaluations at the same point agree exactly and it can be differentiated,
for example when training through it.

\subsection{What Determines the Accuracy}
\label{sec:accuracy-pcpg}

\subsubsection{The scale of $z$}

\subsubsection{The missing ratio}

\subsubsection{The number of features}


\section{Training with PCPG}
\label{sec:disc-training}

\subsection{Why the estimate looks better but accuracy doesn't}

\subsection{Why MC isn't much worse}


\subsection{Why isn't standardised data better in accuracy}

\subsection{Takeaway}

\section{Related Work}

Here, we just focus on works that also try to calculate 

\[
  \mathbb{E}_{P(\xm \mid \xo)} \left[\sigma(w^T x)\right]
\]

but use a different approach.

\subsection{Naive Conformant Learning (NaCL)}

\citet{khosravi2019expect} first point out that this expectation is hard
to compute in general: even with a naive Bayes feature distribution, it
is NP-hard. They also note that the popular shortcut of mean imputation
only matches the expectation when the classifier is linear \emph{and} the
features are independent --- for the sigmoid it is always biased.

Their fix is to pick a feature distribution on which the expectation
\emph{is} cheap, namely a naive Bayes model over binary features. The
trick is that naive Bayes and logistic regression are dual to each other:
every naive Bayes induces a logistic regression, and many naive Bayes
models induce the same one. They call a naive Bayes \textit{conformant}
with the given LR if both give the same class probabilities on fully
observed inputs, and learn such a distribution via geometric programming
(``naive conformant learning''). Once this naive Bayes is fitted, the
expected prediction is just a marginal inference query and runs in linear
time.

\medskip
\noindent\textbf{Difference to us:} they keep the feature distribution
simple (naive Bayes on binary inputs) so the expectation is exact and
cheap; we keep the feature distribution rich (a GMM or PC on continuous
inputs) and instead rewrite the sigmoid with two identities. They
also have to refit the feature distribution whenever the classifier
changes, while ours is trained once and reused.

\subsection{Joint Modelling via SAEM}

\citet{jiang2020logreg} assume a joint Gaussian model on the covariates,
$x \sim \mathcal{N}(\mu, \Sigma)$, together with logistic regression on
top. Their main contribution is on the training side: they fit
$(\mu, \Sigma, \beta)$ jointly from training data that may itself contain
missing values, using a stochastic-approximation EM algorithm with
Metropolis--Hastings updates.

For prediction at test time, they simply do plain Monte Carlo: sample
$x^{m,(s)} \sim p(\xm \mid \xo)$ from the fitted Gaussian and average,
\[
  p(y \mid \xo) \approx \frac{1}{S} \sum_{s=1}^{S} p\!\left(y \mid \xo, x^{m,(s)}\right)
\]
This is exactly the MC baseline we compare against, just with a single
Gaussian instead of a mixture as the density.

\medskip
\noindent\textbf{Difference to us:} they sample the missing features and
average the sigmoid; we never sample the missing features, but instead
rewrite the sigmoid so that only the characteristic function of
$p(\xm \mid \xo)$ is needed. Their feature distribution is also limited
to a single Gaussian, while ours can capture multimodal shapes.
